\documentclass[11pt,a4paper]{article}

% ── Packages ──────────────────────────────────────────────────────────────────
\usepackage{amsmath,amssymb}
\usepackage{booktabs}
\usepackage{graphicx}
\usepackage{float}
\usepackage{adjustbox}
\usepackage[numbers,sort&compress]{natbib}
\usepackage{hyperref}
\usepackage{geometry}
\usepackage{microtype}
\usepackage{xcolor}

\geometry{margin=2.5cm}

% ── Title block ───────────────────────────────────────────────────────────────
\title{KGE-Based Diagnostic Framework for Variance Collapse Detection\\
       in Multi-Horizon PM$_{10}$ Forecasting}

\author{Francisco Federico Garcia Crespi\\
        \small Universidad Miguel Hern\'andez de Elche\\
        \small \texttt{fedeg@umh.es}}

\date{}

% ──────────────────────────────────────────────────────────────────────────────
\begin{document}
\maketitle

% ── Abstract ──────────────────────────────────────────────────────────────────
\begin{abstract}
Multi-horizon forecasting of particulate matter (PM$_{10}$) concentrations
is a critical task for air-quality management, yet conventional accuracy
metrics such as RMSE fail to diagnose \emph{variance collapse} — the
systematic underdispersion that afflicts many machine-learning models at
longer lead times.
%
We propose a diagnostic framework grounded in the Kling-Gupta Efficiency
(KGE) decomposition~\citep{gupta2009} that separates forecast error into
three interpretable components: correlation ($r$), variability ($\alpha$),
and bias ($\beta$).
%
We introduce \emph{Variance Retention} (VR) as an explicit diagnostic, and
combine it with the RMSE-based skill score into a composite diagnostic
indicator (Skill$_{\mathrm{VP}}$) that penalises skill claims made by
variance-collapsed models.
%
Experiments are conducted on hourly PM$_{10}$ records from three EEA
monitoring stations in Madrid (2019–2022) under a strictly leakage-free
rolling-origin evaluation protocol.
%
Four forecasting strategies are benchmarked: persistence, ARIMA, LSTM, and
LightGBM, at horizons $h \in \{1, 6, 12, 24\}$~h.
%
Results show that \ldots [TO BE COMPLETED].

\medskip
\noindent\textbf{Keywords:} air quality forecasting, PM$_{10}$, KGE, variance
collapse, rolling-origin cross-validation, multi-horizon forecasting,
machine learning.
\end{abstract}

% ── 1. Introduction ───────────────────────────────────────────────────────────
\section{Introduction}\label{sec:intro}

Air pollution forecasting is a central tool for environmental health
management~\citep{WHO2021}.
Hourly predictions of PM$_{10}$ concentrations at 1-h to 24-h lead times
directly inform short-term advisory systems, yet reliable evaluation
remains methodologically challenging.

Evaluation frameworks for time-series forecasting typically rely on RMSE or
MAE, which measure average accuracy but are insensitive to the statistical
spread of predictions.
A model that always predicts the climatological mean may achieve competitive
RMSE while exhibiting zero predictive variance — a pathology known as
\emph{variance collapse}~\citep{Gneiting2007}.
This failure mode is particularly common in neural-network and
gradient-boosting regressors at longer horizons.

The Kling-Gupta Efficiency (KGE)~\citep{gupta2009}, originally developed for
hydrological streamflow forecasting, decomposes the forecast-observation
relationship into three orthogonal components.
Unlike RMSE, KGE penalises models whose variability ($\alpha$ component)
deviates from 1.0, making it inherently sensitive to variance collapse.

\paragraph{Contributions.}
\begin{enumerate}
  \item A diagnostic framework that applies KGE decomposition to multi-horizon
        PM$_{10}$ forecasting to detect variance collapse (Section~\ref{sec:method}).
  \item A formal definition of Variance Retention (VR) as a standalone
        diagnostic (Section~\ref{sec:vr}).
  \item Skill$_{\mathrm{VP}}$: a composite diagnostic that couples RMSE skill
        and VR (Section~\ref{sec:skillvp}).
  \item Empirical evaluation on Madrid EEA data under a leakage-free rolling-
        origin protocol (Section~\ref{sec:exp}).
\end{enumerate}

% ── 2. Related Work ───────────────────────────────────────────────────────────
\section{Related Work}\label{sec:related}

\subsection{Multi-horizon air quality forecasting}

[TODO: review of PM10 forecasting literature — ARIMA, gradient boosting, LSTM]

\subsection{Evaluation of probabilistic and deterministic forecasts}

[TODO: proper scoring rules, CRPS, NRG, skill scores, variance diagnostics]

\subsection{Rolling-origin evaluation}

[TODO: Tashman 2000, Bergmeir 2018 on leakage in time-series cross-validation]

\subsection{KGE in hydrology and beyond}

[TODO: Gupta 2009, Knoben 2019, applications outside hydrology]

% ── 3. Methodology ────────────────────────────────────────────────────────────
\section{Methodology}\label{sec:method}

\subsection{Rolling-Origin Evaluation Protocol}\label{sec:protocol}

Let $\{x_t\}_{t=1}^{T}$ denote the full hourly PM$_{10}$ series.
We partition it into $K=16$ rolling-origin folds as follows.
The $k$-th fold ($k=0,\ldots,K-1$) uses a training window
$\mathcal{T}_k = \{x_t : t \le n_k\}$ of size
$n_k = n_{\min} + k \cdot \lfloor(T - n_{\min} - h_{\max})/K\rfloor$,
where $n_{\min} = 8760$ (one year of hourly data).
For each horizon $h \in \mathcal{H} = \{1, 6, 12, 24\}$, the test target is
$x_{n_k + h}$, which is strictly outside $\mathcal{T}_k$.

\emph{Leakage-free preprocessing}: all data transformations (scaling,
imputation) are fitted exclusively on $\mathcal{T}_k$ and applied to the
test target using the training-fitted parameters.
No information from future time points enters any fold's training stage.

\subsection{KGE Decomposition}\label{sec:kge}

The Kling-Gupta Efficiency~\citep{gupta2009} is defined as:
\begin{equation}
  \mathrm{KGE} = 1 - \sqrt{(r-1)^2 + (\alpha-1)^2 + (\beta-1)^2},
  \label{eq:kge}
\end{equation}
where, given observed values $\mathbf{o}$ and forecast values $\mathbf{f}$:
\begin{align}
  r      &= \operatorname{Corr}(\mathbf{o},\mathbf{f}), \label{eq:r}\\
  \alpha &= \frac{\sigma_f}{\sigma_o}, \label{eq:alpha}\\
  \beta  &= \frac{\mu_f}{\mu_o}. \label{eq:beta}
\end{align}
A perfect forecast achieves $r = \alpha = \beta = 1$ and
$\mathrm{KGE} = 1$.
Values below $-0.41$ indicate performance worse than the climatological
mean~\citep{Knoben2019}.

\subsection{Variance Retention}\label{sec:vr}

The variability component $\alpha$ in Eq.~\eqref{eq:alpha} is equivalent
to the square root of the ratio of predicted to observed variance.
We define \emph{Variance Retention} (VR) as:
\begin{equation}
  \mathrm{VR} = \frac{\operatorname{Var}(\mathbf{f})}{\operatorname{Var}(\mathbf{o})}
               = \alpha^2.
  \label{eq:vr}
\end{equation}
VR$<1$ indicates underdispersion; VR$>1$ indicates overdispersion.
A model suffering complete variance collapse has VR$\approx 0$.

\subsection{RMSE Skill Score}\label{sec:skill}

\begin{equation}
  \mathrm{Skill}_{\mathrm{RMSE}} = 1 - \frac{\mathrm{RMSE}_{\mathrm{model}}}
                                            {\mathrm{RMSE}_{\mathrm{persistence}}}.
  \label{eq:skill}
\end{equation}

Positive values indicate improvement over persistence.

\subsection{Skill$_{\mathrm{VP}}$ Diagnostic}\label{sec:skillvp}

A model may report high Skill$_{\mathrm{RMSE}}$ yet exhibit severe variance
collapse, which is undetectable by the skill score alone.
We define the \emph{Variance-Penalised Skill} as:
\begin{equation}
  \mathrm{Skill}_{\mathrm{VP}} = \mathrm{Skill}_{\mathrm{RMSE}}
                                 \cdot \min(1,\,\mathrm{VR}).
  \label{eq:skillvp}
\end{equation}
When VR$\ge 1$ the penalty is inactive.
When VR$<1$ the skill is deflated proportionally to the degree of
underdispersion.
This is a \emph{diagnostic} composite rather than a standalone metric,
designed to reveal discrepancies between accuracy and calibration.

% ── 4. Experimental Setup ─────────────────────────────────────────────────────
\section{Experimental Setup}\label{sec:exp}

\subsection{Data}

Hourly PM$_{10}$ measurements for the period 2019–2022 were obtained from
the European Environment Agency (EEA) e-Reporting database for three Madrid
monitoring stations: ES1422A (Plaza de Espa\~na), ES1938A (Escuelas
Aguirre), and ES1942A (Arturo Soria).
Coverage statistics are reported in Table~\ref{tab:data}.

\begin{table}[H]
  \centering
  \caption{Dataset coverage statistics.}
  \label{tab:data}
  % \input{../results/tables/data_coverage.tex}
  [PLACEHOLDER — replace with \texttt{\textbackslash input\{results/tables/data\_coverage.tex\}}]
\end{table}

\subsection{Models}

[TODO: brief description of each model and hyperparameters]

\subsection{Evaluation}

[TODO: horizon profile computation, averaging over folds and stations]

% ── 5. Results ────────────────────────────────────────────────────────────────
\section{Results}\label{sec:results}

\subsection{Accuracy profiles}

\begin{table}[H]
  \centering
  \caption{Full metric comparison averaged over 16 folds and 3 stations.}
  \label{tab:full}
  % \input{../results/tables/full_comparison_latex.tex}
  [PLACEHOLDER]
\end{table}

\begin{figure}[H]
  \centering
  % \includegraphics[width=0.8\linewidth]{../results/figures/skill_rmse.pdf}
  \caption{Skill$_{\mathrm{RMSE}}$ profiles by forecast horizon.}
  \label{fig:skill_rmse}
\end{figure}

\subsection{Variance collapse analysis}

\begin{figure}[H]
  \centering
  % \includegraphics[width=0.9\linewidth]{../results/figures/variance_retention.pdf}
  \caption{Variance Retention (VR) by forecast horizon.
           Dashed line at VR$=1$ marks the collapse threshold.}
  \label{fig:vr}
\end{figure}

\begin{figure}[H]
  \centering
  % \includegraphics[width=\linewidth]{../results/figures/kge_components.pdf}
  \caption{KGE component profiles: $r$ (correlation), $\alpha$ (variability),
           $\beta$ (bias) by forecast horizon.}
  \label{fig:kge_comp}
\end{figure}

\subsection{Skill$_{\mathrm{VP}}$ diagnostic}

\begin{figure}[H]
  \centering
  % \includegraphics[width=0.9\linewidth]{../results/figures/skill_vp.pdf}
  \caption{Skill$_{\mathrm{RMSE}}$ vs Skill$_{\mathrm{VP}}$.
           Divergence between panels signals variance collapse.}
  \label{fig:skill_vp}
\end{figure}

% ── 6. Discussion ─────────────────────────────────────────────────────────────
\section{Discussion}\label{sec:discussion}

[TODO: interpret rank reversals between Skill$_{\mathrm{RMSE}}$ and
Skill$_{\mathrm{VP}}$; discuss practical implications for model selection;
limitations of the framework]

% ── 7. Conclusion ─────────────────────────────────────────────────────────────
\section{Conclusion}\label{sec:conclusion}

[TODO]

% ── References ────────────────────────────────────────────────────────────────
\bibliographystyle{elsarticle-num}
\bibliography{references}

\end{document}
